\chapter*{Introduction}
\addcontentsline{toc}{chapter}{Introduction}
\label{ch:introduction}
\index{startup sequence}\index{garbage collection}\index{keybinding dispatch}\index{LSP!integration}\index{org-mode!export pipeline}

This book is a deep tour of the mechanisms that make GNU~Emacs
tick: from the two-file startup sequence and garbage-collection
tuning to the keybinding dispatch stack, the LSP integration layer,
and the org-mode export pipeline.

\paragraph{Who this book is for.}
You have been using Emacs for a while.  You can install packages,
write basic Elisp, and navigate the Info manual.  What you lack is
the mental model that ties the pieces together: \emph{why} does
\texttt{use-package-always-defer} change loading behaviour?
\emph{How} does \texttt{eglot-ensure} select a server?  \emph{When}
does \texttt{gc-cons-threshold} actually matter?  This book answers
those questions.

\paragraph{How to read it.}
The six parts are largely independent.  Part~\ref{part:startup}
(startup and performance) is foundational and worth reading first;
Part~\ref{part:packages} (the loading model) is a close second,
since its concepts surface in every later chapter.  After those two,
jump to whichever part addresses your current interest.  Each chapter
opens with a signpost that sets expectations and closes with a
\textbf{Mental model} summary you can use as a revision checklist.

\paragraph{Expert-depth blocks.}
Paragraphs marked with a blue \textit{Expert depth} sidebar contain
internal implementation detail that is not needed on a first reading.
Skip them unless you are debugging a specific problem or are simply
curious.

\paragraph{Appendices.}
The appendices collect quick-reference tables (variables, hooks,
keybindings) and deep-dive material extracted from the main chapters
so that the narrative flow is not interrupted by long reference
listings.
